\documentclass[11pt]{ctexart}
\usepackage[a4paper,margin=1in]{geometry}
\usepackage{longtable,booktabs}
\title{Geant4-Agent Casebank Regression Report (EN)}
\author{Automated Pipeline}
\date{2026-03-08\_102153}
\begin{document}
\maketitle
\section{Unified Pass/Fail Table}
\begin{longtable}{p{0.06\linewidth}p{0.08\linewidth}p{0.07\linewidth}p{0.08\linewidth}p{0.07\linewidth}p{0.07\linewidth}p{0.07\linewidth}p{0.07\linewidth}p{0.09\linewidth}p{0.06\linewidth}p{0.21\linewidth}}
\toprule
ID & Domain & Style & Completeness & ExpInit & ActInit & ExpFinal & ActFinal & Missing(I->F) & Result & Fail Reason \\
\midrule
MTM01 & medical & fuzzy & missing & FAIL & FAIL & PASS & PASS & 12->0 & PASS & - \\
MTM02 & aerospace & precise & missing & FAIL & FAIL & PASS & PASS & 2->0 & PASS & - \\
MTM03 & industrial\_ndt & precise & missing & FAIL & FAIL & PASS & PASS & 12->0 & PASS & - \\
MTM04 & physics & precise & missing & FAIL & FAIL & PASS & PASS & 8->0 & PASS & - \\
MTM05 & nuclear\_engineering & fuzzy & missing & FAIL & FAIL & PASS & PASS & 5->0 & PASS & - \\
MTM06 & security\_screening & fuzzy & missing & FAIL & FAIL & PASS & PASS & 11->0 & PASS & - \\
MTM07 & semiconductor & precise & missing & FAIL & FAIL & PASS & PASS & 10->0 & PASS & - \\
MTM08 & education & fuzzy & missing & FAIL & FAIL & PASS & PASS & 11->0 & PASS & - \\
MTM09 & environmental\_radiation & fuzzy & missing & FAIL & FAIL & PASS & PASS & 3->0 & PASS & - \\
MTM10 & high\_energy\_physics & precise & missing & FAIL & FAIL & PASS & PASS & 10->0 & PASS & - \\
MTM11 & materials\_update & precise & full & PASS & PASS & PASS & PASS & 0->0 & PASS & - \\
MTM12 & overwrite\_reject & precise & full & PASS & PASS & PASS & PASS & 0->0 & PASS & - \\
MTM13 & medical\_detector & fuzzy & missing & FAIL & FAIL & PASS & PASS & 13->0 & PASS & - \\
MTM14 & industrial\_panel & precise & missing & FAIL & FAIL & PASS & PASS & 15->0 & FAIL & expected\_final\_mismatch:materials.selected\_materials.0 \\
MTM15 & shielding\_stack & precise & missing & FAIL & FAIL & PASS & PASS & 14->0 & PASS & - \\
MTM16 & nested\_target & precise & missing & FAIL & FAIL & PASS & PASS & 13->0 & PASS & - \\
MTM17 & coaxial\_shield & fuzzy & missing & FAIL & FAIL & PASS & PASS & 14->0 & PASS & - \\
MTM18 & boolean\_cavity & precise & missing & FAIL & FAIL & PASS & PASS & 17->0 & FAIL & expected\_final\_mismatch:materials.selected\_materials.0 \\
MTM19 & ring\_complete & precise & full & PASS & PASS & PASS & PASS & 0->0 & PASS & - \\
MTM20 & grid\_complete & precise & full & PASS & FAIL & PASS & FAIL & 1->1 & FAIL & initial\_complete\_mismatch, final\_complete\_mismatch, expected\_final\_mismatch:materials.selected\_materials.0 \\
MTM21 & layered\_detector & fuzzy & missing & FAIL & FAIL & PASS & PASS & 14->0 & PASS & - \\
MTM22 & embedded\_sensor & precise & full & PASS & PASS & PASS & PASS & 0->0 & PASS & - \\
MTM23 & layered\_shield\_complete & precise & full & PASS & PASS & PASS & PASS & 0->0 & PASS & - \\
MTM24 & boolean\_complete & precise & full & PASS & FAIL & PASS & FAIL & 1->1 & FAIL & initial\_complete\_mismatch, final\_complete\_mismatch, expected\_final\_mismatch:materials.selected\_materials.0, expected\_final\_mismatch:source.type \\
\bottomrule
\end{longtable}
\section{Baseline Comparison}
\begin{itemize}
\item pass\_rate=0.833
\item initial\_alignment\_rate=0.917
\item final\_alignment\_rate=0.917
\item avg\_turns\_per\_case=2.417
\item missing\_reduction\_rate=1.000
\item confirm\_apply\_success\_rate=1.000
\item rollback\_after\_confirm\_rate=0.000
\item internal\_leak\_turn\_rate=0.034
\item repeated\_assistant\_case\_rate=0.000
\item stale\_confirm\_case\_rate=0.000
\item delta\_pass\_rate=-0.117
\item regressions=none
\item improvements=none
\end{itemize}
\appendix
\section{Appendix: Full Dialogues and Naturalized Internal Chains}
\subsection{Case MTM01 (medical)}
\textbf{Dialogue:}
\begin{itemize}
\item User: 我先要一个用于医学成像的初步辐照场景，细节还没定，你先告诉我还缺什么。
\item Assistant: 好的，我们来看看目前的配置进度。

**已完成：**
*   几何部分：已经确定了使用“盒体”作为基础结构。

**待补充：**
*   几何：需要您提供盒体在X、Y、Z三个方向的具体尺寸。
*   材料：需要为这个盒体体积指定对应的材料。
*   源：需要确定放射源的类型（例如伽马源）、粒子种类和能量。
*   物理：需要选择模拟中使用的物理过程列表。
*   输出：需要设定输出数据的格式和保存路径。

请告诉我盒体的尺寸，我们继续往下配置。
\item User: 采用 300 mm x 300 mm x 300 mm 的水箱，材料 ?；源用束流 gamma，能量 6 MeV，位置 (0,0,-400) mm，方向 +z；物理列表 QBBC；输出 root。
\item Assistant: 我注意到这可能会覆盖已经确认的设置。具体来说：几何结构类型将从“未知”改为“单盒体”；盒体 x 方向尺寸将从“未设置”改为“300.0”。如果您确认应用这些修改，请直接回复“确认”；如果您希望保留原来的值，请直接回复“保持原值”。
\item User: 确认
\item Assistant: 配置已完成！目前几何部分的盒体尺寸和类型都已确认。
\end{itemize}
\textbf{Internal Chain (Natural Language):}
\begin{itemize}
\item Turn 1: action=summarize\_progress, phase=geometry, complete=False.
\item Missing fields: geometry.params.module\_x, geometry.params.module\_y, geometry.params.module\_z, materials.volume\_material\_map, source.type, source.particle, source.energy, source.position, source.direction, physics.physics\_list, output.format, output.path.
\item LLM used=True; fallback=none; backend=llm\_semantic\_frame+runtime\_semantic.
\item Turn 2: action=confirm\_overwrite, phase=geometry, complete=False.
\item Missing fields: geometry.params.module\_x, geometry.params.module\_y, geometry.params.module\_z.
\item LLM used=True; fallback=none; backend=llm\_slot\_frame+runtime\_semantic.
\item Turn 3: action=finalize, phase=materials, complete=True.
\item Missing fields: none.
\item LLM used=True; fallback=none; backend=llm\_slot\_frame+runtime\_semantic+confirmed\_pending\_overwrite.
\end{itemize}
\subsection{Case MTM02 (aerospace)}
\textbf{Dialogue:}
\begin{itemize}
\item User: I need a first-pass proton irradiation setup for a silicon device; keep it simple for now.
\item Assistant: Great progress so far. Your geometry, materials, source, and physics settings are now configured.

To finish, I still need to know:
*   Your preferred output format.
*   The file path where you'd like to save the output.
\item User: Use a silicon box 50 mm x 50 mm x 5 mm, beam source proton 100 MeV from (0,0,-100) mm along +z, physics FTFP\_BERT, output json.
\item Assistant: You've already set a box size. This change would update it:
*   box size along x: 1000.0 → 50.0
*   box size along y: 1000.0 → 50.0

Please reply **'confirm'** to apply this change, or **'keep existing'** to keep the current value.
\item User: confirm
\item Assistant: Configuration complete. The following items are confirmed: box size along x, box size along y, box size along z, and output path.
\end{itemize}
\textbf{Internal Chain (Natural Language):}
\begin{itemize}
\item Turn 1: action=summarize\_progress, phase=materials, complete=False.
\item Missing fields: output.format, output.path.
\item LLM used=True; fallback=none; backend=llm\_slot\_frame+runtime\_semantic.
\item Turn 2: action=confirm\_overwrite, phase=source, complete=False.
\item Missing fields: none.
\item LLM used=True; fallback=none; backend=llm\_slot\_frame+runtime\_semantic.
\item Turn 3: action=finalize, phase=physics, complete=True.
\item Missing fields: none.
\item LLM used=True; fallback=none; backend=llm\_semantic\_frame+runtime\_semantic+confirmed\_pending\_overwrite.
\end{itemize}
\subsection{Case MTM03 (industrial\_ndt)}
\textbf{Dialogue:}
\begin{itemize}
\item User: 我要做一个 gamma 穿透铁靶的小测试，但我还没把完整参数说完。
\item Assistant: 目前我们已经确定了使用铁作为材料，并选择了 gamma 粒子。接下来还需要你提供几何结构的具体类型和尺寸（比如盒体的长、宽、高），以及指定哪个体积绑定到铁材料。
\item User: 请设为半径 40 mm、半长 80 mm 的铁圆柱，材料 ?；点源 gamma，能量 2 MeV，位置 (0,0,-200) mm，方向 +z；物理列表 Shielding；输出 csv。
\item Assistant: 配置已完成。当前已确认的设置有：体积与材料绑定、输出路径、输出格式和物理列表。
\end{itemize}
\textbf{Internal Chain (Natural Language):}
\begin{itemize}
\item Turn 1: action=summarize\_progress, phase=geometry, complete=False.
\item Missing fields: geometry.structure, geometry.params.module\_x, geometry.params.module\_y, geometry.params.module\_z, materials.volume\_material\_map, source.type, source.energy, source.position, source.direction, physics.physics\_list, output.format, output.path.
\item LLM used=True; fallback=none; backend=llm\_slot\_frame+runtime\_semantic.
\item Turn 2: action=finalize, phase=materials, complete=True.
\item Missing fields: none.
\item LLM used=True; fallback=none; backend=llm\_slot\_frame+runtime\_semantic.
\end{itemize}
\subsection{Case MTM04 (physics)}
\textbf{Dialogue:}
\begin{itemize}
\item User: Gamma point source 1 MeV at (0,0,-100) toward +z; please build the rest.
\item Assistant: Here's a summary of your setup so far and what's still needed.

**Completed:**
*   **Source:** You've configured a gamma point source with 1 MeV energy, positioned at (0,0,-100) and directed along +z.

**Still Required:**
*   **Geometry:** The structure (like a simple box) and its dimensions (size along x, y, and z).
*   **Materials:** You'll need to assign a material to the geometry volume.
*   **Physics \& Output:** A physics list, output format, and file path are also pending.

Please provide the geometry details next.
\item User: Set copper box 1 m x 1 m x 1 m, physics FTFP\_BERT, output hdf5.
\item Assistant: Your configuration is complete. The following items are confirmed: volume-to-material binding, output path, output format, and physics list.
\end{itemize}
\textbf{Internal Chain (Natural Language):}
\begin{itemize}
\item Turn 1: action=summarize\_progress, phase=geometry, complete=False.
\item Missing fields: geometry.structure, geometry.params.module\_x, geometry.params.module\_y, geometry.params.module\_z, materials.volume\_material\_map, physics.physics\_list, output.format, output.path.
\item LLM used=True; fallback=none; backend=llm\_slot\_frame+runtime\_semantic.
\item Turn 2: action=finalize, phase=materials, complete=True.
\item Missing fields: none.
\item LLM used=True; fallback=none; backend=llm\_slot\_frame+runtime\_semantic.
\end{itemize}
\subsection{Case MTM05 (nuclear\_engineering)}
\textbf{Dialogue:}
\begin{itemize}
\item User: 我想做一个中子屏蔽的概念验证，现在先给一个缺参开头。
\item Assistant: 目前我们已经完成了几何结构、材料、源类型、物理列表和输出设置的基本同步。

接下来还需要补充以下参数才能继续：
- 几何：盒体在x、y、z三个方向的尺寸
- 源：粒子的能量大小和发射位置

请提供这些信息，我们就可以开始配置了。
\item User: 使用 200 mm x 200 mm x 200 mm 的铅块，材料 ?；点源 neutron，能量 14 MeV，位置 (0,0,-50) mm，方向 +z；物理列表 QGSP\_BERT\_HP；输出 xml。
\item Assistant: 我检测到这会覆盖已经确认的值。待修改内容：几何结构类型：unknown -> single\_box；已选材料：['G4\_Cu'] -> ['G4\_Pb']。如果要应用这些修改，请直接回复“确认”；如果要保留原值，请直接说“保持原值”。
\item User: 确认
\item Assistant: 配置已完成。当前已确认的参数包括：体积与材料绑定、输出路径，以及盒体在 z、y、x 三个方向的尺寸。
\end{itemize}
\textbf{Internal Chain (Natural Language):}
\begin{itemize}
\item Turn 1: action=summarize\_progress, phase=geometry, complete=False.
\item Missing fields: geometry.params.module\_x, geometry.params.module\_y, geometry.params.module\_z, source.energy, source.position.
\item LLM used=True; fallback=none; backend=llm\_slot\_frame+runtime\_semantic.
\item Turn 2: action=confirm\_overwrite, phase=geometry, complete=False.
\item Missing fields: geometry.params.module\_x, geometry.params.module\_y, geometry.params.module\_z.
\item LLM used=True; fallback=none; backend=llm\_slot\_frame+runtime\_semantic.
\item Turn 3: action=finalize, phase=materials, complete=True.
\item Missing fields: none.
\item LLM used=True; fallback=none; backend=llm\_slot\_frame+runtime\_semantic+confirmed\_pending\_overwrite.
\end{itemize}
\subsection{Case MTM06 (security\_screening)}
\textbf{Dialogue:}
\begin{itemize}
\item User: Need a quick cargo-screening toy setup with forward photons, but I have not fixed the target shape, material, source energy, or output format yet.
\item Assistant: Here's what we have so far:
*   **Geometry:** You've chosen the structure. Next, I need the box dimensions (X, Y, and Z).
*   **Materials:** I need to know what material to assign to your volume.
*   **Source:** You've selected photons. Next, I need the source type and its energy.
*   **Physics:** I need the physics list.
*   **Output:** I need the output format and file path.
\item User: Use steel cylinder target radius 60 mm half-length 120 mm, beam gamma 3 MeV from (0,0,-500) mm to +z, physics Shielding, output root.
\item Assistant: This would overwrite an already confirmed value. Pending change: geometry structure: unknown -> single\_tubs; box size along x: None -> 1000.0. Reply 'confirm' to apply it, or say 'keep existing' to keep the current value.
\item User: confirm
\item Assistant: Configuration complete. The following items are now confirmed: the volume-to-material binding, the half length, the radius, and the box size along Z.
\end{itemize}
\textbf{Internal Chain (Natural Language):}
\begin{itemize}
\item Turn 1: action=summarize\_progress, phase=geometry, complete=False.
\item Missing fields: geometry.params.module\_x, geometry.params.module\_y, geometry.params.module\_z, materials.volume\_material\_map, source.type, source.energy, source.position, source.direction, physics.physics\_list, output.format, output.path.
\item LLM used=True; fallback=none; backend=llm\_slot\_frame+runtime\_semantic.
\item Turn 2: action=confirm\_overwrite, phase=geometry, complete=False.
\item Missing fields: geometry.params.module\_x, geometry.params.module\_y, geometry.params.module\_z.
\item LLM used=True; fallback=none; backend=llm\_slot\_frame+runtime\_semantic.
\item Turn 3: action=finalize, phase=materials, complete=True.
\item Missing fields: none.
\item LLM used=True; fallback=none; backend=llm\_semantic\_frame+runtime\_semantic+confirmed\_pending\_overwrite.
\end{itemize}
\subsection{Case MTM07 (semiconductor)}
\textbf{Dialogue:}
\begin{itemize}
\item User: 我想搭一个硅薄片的中子照射测试，但先只给一个简述。
\item Assistant: 好的，目前我们已经确认了以下几个部分：

*   **几何**：结构确定为盒体，但还需要您提供长、宽、高的具体尺寸。
*   **材料**：已选定为硅，并会将其与几何体绑定。
*   **源**：粒子类型已确定为中子，接下来需要您指定源的种类（如点源、面源等）和能量。

为了完成配置，接下来还需要您补充物理过程列表以及输出相关的设置。请随时告诉我这些信息。
\item User: 使用 20 mm x 20 mm x 1 mm 的硅片，材料 ?；各向同性 neutron，能量 2 MeV，位置 (0,0,0) mm；物理列表 QGSP\_BERT\_HP；输出 json。
\item Assistant: 配置已完成。当前已确认的设置包括：输出路径、输出格式、物理列表和源方向。
\end{itemize}
\textbf{Internal Chain (Natural Language):}
\begin{itemize}
\item Turn 1: action=summarize\_progress, phase=geometry, complete=False.
\item Missing fields: geometry.params.module\_x, geometry.params.module\_y, geometry.params.module\_z, source.type, source.energy, source.position, source.direction, physics.physics\_list, output.format, output.path.
\item LLM used=True; fallback=none; backend=llm\_slot\_frame+runtime\_semantic.
\item Turn 2: action=finalize, phase=materials, complete=True.
\item Missing fields: none.
\item LLM used=True; fallback=none; backend=llm\_slot\_frame+runtime\_semantic.
\end{itemize}
\subsection{Case MTM08 (education)}
\textbf{Dialogue:}
\begin{itemize}
\item User: For classroom demo, I want gamma attenuation in copper; you can ask only essential missing parts.
\item Assistant: Here's a summary of your setup so far and what's still needed:

**Completed:**
*   **Materials:** Copper is selected.
*   **Source:** Gamma particle is set.
*   **Physics:** A physics list has been chosen.

**Still needed to configure:**
*   **Geometry:** The structure type and the box size along the X, Y, and Z axes.
*   **Materials:** Assigning the copper material to a specific volume.
*   **Source:** The source type, its energy, and position.
*   **Output:** The output format and file path.
\item User: Okay, copper box 100 mm x 100 mm x 100 mm, point gamma 1 MeV at (0,0,-200) mm, direction +z, physics FTFP\_BERT, output csv.
\item Assistant: You're about to overwrite a confirmed setting.

**Pending change:**  
Physics list: `Shielding` → `FTFP\_BERT`

Please reply **'confirm'** to apply this change, or **'keep existing'** to keep the current value.
\item User: confirm
\item Assistant: Configuration complete. The following items are now confirmed: physics list, output path, output format, and source direction.
\end{itemize}
\textbf{Internal Chain (Natural Language):}
\begin{itemize}
\item Turn 1: action=summarize\_progress, phase=geometry, complete=False.
\item Missing fields: geometry.structure, geometry.params.module\_x, geometry.params.module\_y, geometry.params.module\_z, materials.volume\_material\_map, source.type, source.energy, source.position, source.direction, output.format, output.path.
\item LLM used=True; fallback=none; backend=llm\_slot\_frame+runtime\_semantic.
\item Turn 2: action=confirm\_overwrite, phase=materials, complete=False.
\item Missing fields: none.
\item LLM used=True; fallback=none; backend=llm\_slot\_frame+runtime\_semantic.
\item Turn 3: action=finalize, phase=source, complete=True.
\item Missing fields: none.
\item LLM used=True; fallback=none; backend=llm\_semantic\_frame+runtime\_semantic+confirmed\_pending\_overwrite.
\end{itemize}
\subsection{Case MTM09 (environmental\_radiation)}
\textbf{Dialogue:}
\begin{itemize}
\item User: 我想先搭一个空气体积中的低能 gamma 传播场景，细节稍后补。
\item Assistant: 目前我们已经完成了几何结构、材料和源的设置，接下来还需要你补充物理列表、输出格式和输出路径这三项信息。
\item User: 设为 500 mm x 500 mm x 500 mm 的空气盒，材料 ??；点源 gamma，能量 0.5 MeV，位置 (0,0,-300) mm，方向 +z；物理列表 QBBC；输出 root。
\item Assistant: 我注意到这将会覆盖已经确认的设置。具体修改如下：

- 盒体 x 方向尺寸：1000.0 → 500.0
- 盒体 y 方向尺寸：1000.0 → 500.0

请确认你的选择：
- 如果应用这些修改，请回复“确认”。
- 如果想保留原来的数值，请回复“保持原值”。
\item User: 确认
\item Assistant: 配置已完成。当前已确认的设置包括：源位置、源能量、盒体 z 方向尺寸、盒体 y 方向尺寸和盒体 x 方向尺寸。
\end{itemize}
\textbf{Internal Chain (Natural Language):}
\begin{itemize}
\item Turn 1: action=summarize\_progress, phase=materials, complete=False.
\item Missing fields: physics.physics\_list, output.format, output.path.
\item LLM used=True; fallback=none; backend=llm\_slot\_frame+runtime\_semantic.
\item Turn 2: action=confirm\_overwrite, phase=source, complete=False.
\item Missing fields: none.
\item LLM used=True; fallback=none; backend=llm\_slot\_frame+runtime\_semantic.
\item Turn 3: action=finalize, phase=physics, complete=True.
\item Missing fields: none.
\item LLM used=True; fallback=none; backend=llm\_semantic\_frame+runtime\_semantic+confirmed\_pending\_overwrite.
\end{itemize}
\subsection{Case MTM10 (high\_energy\_physics)}
\textbf{Dialogue:}
\begin{itemize}
\item User: Set geometry as a sphere target; I have not decided radius or source details yet.
\item Assistant: Here's where we stand:

**Geometry:** You've set the structure to a sphere. We still need the sphere's radius.

**Materials:** We need to assign a material to your volume.

**Source:** We need the source type, particle type, and particle energy.

**Physics \& Output:** We also need the physics list, output format, and output path.
\item User: Sphere target radius 80 mm, lead material, electron point source 5 MeV at (0,0,-200) mm toward +z, physics FTFP\_BERT, output hdf5.
\item Assistant: Your configuration is complete. The following items have been confirmed: the volume-to-material binding, output path, output format, and physics list.
\end{itemize}
\textbf{Internal Chain (Natural Language):}
\begin{itemize}
\item Turn 1: action=summarize\_progress, phase=geometry, complete=False.
\item Missing fields: geometry.params.child\_rmax, materials.volume\_material\_map, source.type, source.particle, source.energy, source.position, source.direction, physics.physics\_list, output.format, output.path.
\item LLM used=True; fallback=none; backend=llm\_slot\_frame+runtime\_semantic.
\item Turn 2: action=finalize, phase=materials, complete=True.
\item Missing fields: none.
\item LLM used=True; fallback=none; backend=llm\_slot\_frame+runtime\_semantic.
\end{itemize}
\subsection{Case MTM11 (materials\_update)}
\textbf{Dialogue:}
\begin{itemize}
\item User: Create a complete setup: 100 mm x 100 mm x 100 mm copper box, point gamma 1 MeV at (0,0,-200) mm toward +z, physics QBBC, output json.
\item Assistant: Your configuration is complete. The following items have been confirmed: volume-to-material binding, output path, and output format.
\item User: Change the target material to aluminum and change the output format to root.
\item Assistant: You're about to overwrite two confirmed settings:

*   **Selected materials:** Copper → Aluminum
*   **Output format:** JSON → ROOT

Please reply **'confirm'** to apply these changes, or **'keep existing'** to keep the current values.
\item User: confirm
\item Assistant: Your configuration is complete. The following items have been confirmed: selected materials, output format, output path, and volume-to-material binding.
\end{itemize}
\textbf{Internal Chain (Natural Language):}
\begin{itemize}
\item Turn 1: action=finalize, phase=materials, complete=True.
\item Missing fields: none.
\item LLM used=True; fallback=none; backend=llm\_slot\_frame+runtime\_semantic.
\item Turn 2: action=confirm\_overwrite, phase=source, complete=False.
\item Missing fields: none.
\item LLM used=True; fallback=none; backend=llm\_slot\_frame+runtime\_semantic.
\item Turn 3: action=finalize, phase=physics, complete=True.
\item Missing fields: none.
\item LLM used=True; fallback=none; backend=llm\_semantic\_frame+runtime\_semantic+confirmed\_pending\_overwrite.
\end{itemize}
\subsection{Case MTM12 (overwrite\_reject)}
\textbf{Dialogue:}
\begin{itemize}
\item User: 请创建一个完整配置：100 mm x 100 mm x 100 mm 的铜盒，材料 ?；点源 gamma，能量 1 MeV，位置 (0,0,-200) mm，方向 +z；物理列表 QBBC；输出 json。
\item Assistant: 配置已完成。当前已确认的设置包括：体积与材料绑定、输出路径和输出格式。
\item User: 把材料改成 ?，并把输出格式改成 root。
\item Assistant: 我注意到这次修改会覆盖已确认的设置。具体变更如下：

- 已选材料：从铜（G4\_Cu）改为铝（G4\_Al）
- 输出格式：从 json 改为 root

请确认是否应用这些修改：
• 回复“确认”将应用以上变更
• 回复“保持原值”则保留原有设置
\item User: 保持原值
\item Assistant: 好的，已保留已选材料和输出格式的原有值，这次修改没有应用。
\end{itemize}
\textbf{Internal Chain (Natural Language):}
\begin{itemize}
\item Turn 1: action=finalize, phase=materials, complete=True.
\item Missing fields: none.
\item LLM used=True; fallback=none; backend=llm\_slot\_frame+runtime\_semantic.
\item Turn 2: action=confirm\_overwrite, phase=source, complete=False.
\item Missing fields: geometry.ask.boolean.solid\_a\_size, geometry.ask.boolean.solid\_b\_size.
\item LLM used=True; fallback=none; backend=llm\_slot\_frame+runtime\_semantic.
\item Turn 3: action=reject\_overwrite, phase=physics, complete=True.
\item Missing fields: none.
\item LLM used=True; fallback=none; backend=llm\_semantic\_frame+runtime\_semantic+rejected\_pending\_overwrite.
\end{itemize}
\subsection{Case MTM13 (medical\_detector)}
\textbf{Dialogue:}
\begin{itemize}
\item User: I want a ring of detector modules around a small source region, but I have not fixed the ring radius or module size yet.
\item Assistant: Your geometry structure is set. Next, I need the ring module size and count for the detector ring, and to assign materials to your volumes.

We also need to define your particle source (type, particle, and energy), choose a physics list, and set up your output format and file path.
\item User: Use a ring of 12 modules, each 8 mm x 10 mm x 2 mm, radius 40 mm, clearance 1 mm, material water; point gamma source 0.662 MeV at (0,0,0) mm; physics QBBC; output json.
\item Assistant: You are about to overwrite a confirmed geometry setting.

The pending change is:
- Module size: from 7.83 x 4.72 x 6.86 mm to 8 x 10 x 2 mm.
- Ring: from 4 modules with a 67.49 mm radius and 0.19 mm clearance to 12 modules with a 40 mm radius and 1 mm clearance.

Please reply **'confirm'** to apply this change, or **'keep existing'** to keep the current values.
\item User: confirm
\item Assistant: Configuration complete. The following geometry items are confirmed: ring clearance, ring radius, ring module count, and ring module size.
\end{itemize}
\textbf{Internal Chain (Natural Language):}
\begin{itemize}
\item Turn 1: action=summarize\_progress, phase=materials, complete=False.
\item Missing fields: materials.volume\_material\_map, source.type, source.particle, source.energy, source.position, source.direction, physics.physics\_list, output.format, output.path, geometry.ask.ring.module\_size, geometry.ask.ring.count, geometry.ask.ring.radius, geometry.ask.ring.clearance.
\item LLM used=True; fallback=none; backend=llm\_slot\_frame+runtime\_semantic.
\item Turn 2: action=confirm\_overwrite, phase=source, complete=False.
\item Missing fields: none.
\item LLM used=True; fallback=none; backend=llm\_slot\_frame+runtime\_semantic.
\item Turn 3: action=finalize, phase=physics, complete=True.
\item Missing fields: none.
\item LLM used=True; fallback=none; backend=llm\_semantic\_frame+runtime\_semantic+confirmed\_pending\_overwrite.
\end{itemize}
\subsection{Case MTM14 (industrial\_panel)}
\textbf{Dialogue:}
\begin{itemize}
\item User: 我想做一个二维阵列探测板，但 pitch 和材料还没定。
\item Assistant: 好的，我们已确认了您的几何结构类型（二维阵列探测板）。

接下来，为了完成配置，还需要您提供以下信息：
*   **材料**：体积与材料绑定。
*   **源**：源类型、粒子类型、源能量和源位置。

请随时告诉我这些细节。
\item User: 设成 4 x 3 的盒体阵列，每个模块 10 mm x 10 mm x 2 mm，pitch\_x 12 mm，pitch\_y 14 mm，间隙 1 mm，材料 ?；点源 gamma 0.511 MeV，位置 (0,0,-80) mm，方向 +z；物理列表 QBBC；输出 root。
\item Assistant: 检测到您的修改会覆盖已确认的几何参数。具体来说，阵列模块的尺寸、X/Y方向间距以及阵列间隙都将被更新。

请确认：
- 若**应用此修改**，请直接回复“确认”。
- 若**保留原有数值**，请直接回复“保持原值”。
\item User: 确认
\item Assistant: 好的，配置已完成！目前我们已确认的参数包括：y方向阵列数量、x方向阵列数量、阵列间隙，以及y方向的pitch。
\end{itemize}
\textbf{Internal Chain (Natural Language):}
\begin{itemize}
\item Turn 1: action=summarize\_progress, phase=materials, complete=False.
\item Missing fields: materials.volume\_material\_map, source.type, source.particle, source.energy, source.position, source.direction, physics.physics\_list, output.format, output.path, geometry.ask.grid.module\_size, geometry.ask.grid.count\_x, geometry.ask.grid.count\_y, geometry.ask.grid.pitch\_x, geometry.ask.grid.pitch\_y, geometry.ask.grid.clearance.
\item LLM used=True; fallback=none; backend=llm\_semantic\_frame+runtime\_semantic.
\item Turn 2: action=confirm\_overwrite, phase=source, complete=False.
\item Missing fields: geometry.ask.grid.count\_x, geometry.ask.grid.count\_y.
\item LLM used=True; fallback=none; backend=llm\_slot\_frame+runtime\_semantic.
\item Turn 3: action=finalize, phase=physics, complete=True.
\item Missing fields: none.
\item LLM used=True; fallback=none; backend=llm\_semantic\_frame+runtime\_semantic+confirmed\_pending\_overwrite.
\end{itemize}
\subsection{Case MTM15 (shielding\_stack)}
\textbf{Dialogue:}
\begin{itemize}
\item User: Build a stacked absorber sandwich along z; I will provide the layer sizes next.
\item Assistant: You've set up a stacked absorber geometry along the z-axis. Here's what's complete and what's still needed:

**Completed:**
*   **Geometry:** The basic structure (a stacked absorber sandwich) and the overall box size are set.

**Still Needed:**
*   **Geometry:** The footprint and thickness for each layer of the stack.
*   **Materials:** Which material to assign to each volume.
*   **Source:** The source type, particle type, and its energy.
*   **Physics:** The physics list to use.
*   **Output:** The format and file path for saving results.
\item User: Use stacked layers with footprint 20 mm x 20 mm, thicknesses 2 mm, 4 mm, 6 mm, layer clearance 0.5 mm, outer box 30 mm x 30 mm x 20 mm, container clearance 1 mm, material aluminum; beam gamma 1.25 MeV from (0,0,-120) mm along +z; physics Shielding; output csv.
\item Assistant: This would overwrite an already confirmed geometry configuration.

**Current values:**
*   Stack footprint: \textasciitilde{}16.33 mm x \textasciitilde{}10.28 mm
*   Layer thicknesses: \textasciitilde{}6.68 mm, \textasciitilde{}1.19 mm, \textasciitilde{}5.59 mm
*   Layer clearance: \textasciitilde{}0.73 mm
*   Outer box size: \textasciitilde{}18.07 mm x \textasciitilde{}45.52 mm x \textasciitilde{}16.66 mm
*   Container clearance: \textasciitilde{}0.87 mm

**New values from your request:**
*   Stack footprint: 20.0 mm x 20.0 mm
*   Layer thicknesses: 2.0 mm, 4.0 mm, 6.0 mm
*   Layer clearance: 0.5 mm
*   Outer box size: 30.0 mm x 30.0 mm x 20.0 mm
*   Container clearance: 1.0 mm

Reply **'confirm'** to apply these new values, or **'keep existing'** to keep the current geometry.
\item User: confirm
\item Assistant: Your configuration is now complete. The following geometry items have been confirmed: stack container clearance, stack container size, stack layer clearance, and stack layer thicknesses.
\end{itemize}
\textbf{Internal Chain (Natural Language):}
\begin{itemize}
\item Turn 1: action=summarize\_progress, phase=materials, complete=False.
\item Missing fields: materials.volume\_material\_map, source.type, source.particle, source.energy, source.position, source.direction, physics.physics\_list, output.format, output.path, geometry.ask.stack.footprint, geometry.ask.stack.thicknesses, geometry.ask.stack.layer\_clearance, geometry.ask.stack.parent\_size, geometry.ask.stack.nest\_clearance.
\item LLM used=True; fallback=none; backend=llm\_slot\_frame+runtime\_semantic.
\item Turn 2: action=confirm\_overwrite, phase=source, complete=False.
\item Missing fields: none.
\item LLM used=True; fallback=none; backend=llm\_slot\_frame+runtime\_semantic.
\item Turn 3: action=finalize, phase=physics, complete=True.
\item Missing fields: none.
\item LLM used=True; fallback=none; backend=llm\_semantic\_frame+runtime\_semantic+confirmed\_pending\_overwrite.
\end{itemize}
\subsection{Case MTM16 (nested\_target)}
\textbf{Dialogue:}
\begin{itemize}
\item User: 我想在一个盒子里嵌一个圆柱靶，先别急着补全。
\item Assistant: 好的，我们已确认了“盒体”作为几何结构类型，并记录了其尺寸。

接下来，为了完成模拟配置，我们还需要补充以下信息：
1.  为盒体和内嵌圆柱体分别指定材料。
2.  设置粒子源，包括源的类型、粒子种类和能量。
3.  确定物理过程列表。
4.  选择输出数据的格式和保存路径。

请告诉我您希望如何设置这些部分。
\item User: 外盒 120 mm x 120 mm x 120 mm，内嵌铜圆柱半径 20 mm、半长 40 mm，嵌套间隙 2 mm；材料 ?；束流 gamma 1 MeV，位置 (0,0,-150) mm，方向 +z；物理列表 QBBC；输出 json。
\item Assistant: 我注意到这次修改会覆盖已经确认的几何参数。具体来说，外盒尺寸、内嵌圆柱的半径与半长，以及嵌套间隙都将被更新为新值。

如果您确认要应用这些修改，请直接回复“确认”；如果您希望保留原来的参数，请回复“保持原值”。
\item User: 确认
\item Assistant: 配置已完成。当前已确认的参数包括：嵌套间隙、内嵌圆柱半长、内嵌圆柱半径以及外层容器尺寸。
\end{itemize}
\textbf{Internal Chain (Natural Language):}
\begin{itemize}
\item Turn 1: action=summarize\_progress, phase=materials, complete=False.
\item Missing fields: materials.volume\_material\_map, source.type, source.particle, source.energy, source.position, source.direction, physics.physics\_list, output.format, output.path, geometry.ask.nest.parent\_size, geometry.ask.nest.child\_radius, geometry.ask.nest.child\_half\_length, geometry.ask.nest.clearance.
\item LLM used=True; fallback=none; backend=llm\_slot\_frame+runtime\_semantic.
\item Turn 2: action=confirm\_overwrite, phase=source, complete=False.
\item Missing fields: none.
\item LLM used=True; fallback=none; backend=llm\_slot\_frame+runtime\_semantic.
\item Turn 3: action=finalize, phase=physics, complete=True.
\item Missing fields: none.
\item LLM used=True; fallback=none; backend=llm\_semantic\_frame+runtime\_semantic+confirmed\_pending\_overwrite.
\end{itemize}
\subsection{Case MTM17 (coaxial\_shield)}
\textbf{Dialogue:}
\begin{itemize}
\item User: I need a shell-style shielding toy model, but the shell thicknesses are still open.
\item Assistant: Here's your current setup:

**Geometry:** You've chosen a shell structure. Next, we need the shell's inner radius and thicknesses.
**Materials:** You still need to assign materials to each volume.
**Source:** The particle source type, particle type, and energy are still required.
**Physics:** Your physics list has been set.
**Output:** The output format and file path are still needed.
\item User: Use a shell with inner radius 15 mm, thicknesses 5 mm and 8 mm, half length 40 mm, nested core radius 10 mm, nested core half length 30 mm, clearance 1 mm, material lead; isotropic neutron source 2 MeV at (0,0,0) mm; physics Shielding; output hdf5.
\item Assistant: Your new geometry settings would overwrite the ones you've already confirmed.

**Current geometry:**
*   Shell inner radius: 7.15 mm
*   Shell thicknesses: 1.18 mm, 3.43 mm, 0.86 mm
*   Shell half-length: 22.36 mm
*   Nested core radius: 11.60 mm
*   Nested core half-length: 2.68 mm
*   Clearance: 1.01 mm

**New geometry:**
*   Shell inner radius: 15.0 mm
*   Shell thicknesses: 5.0 mm, 8.0 mm, 0.83 mm
*   Shell half-length: 40.0 mm
*   Nested core radius: 10.0 mm
*   Nested core half-length: 30.0 mm
*   Clearance: 1.0 mm

Please reply **'confirm'** to apply these new values, or **'keep existing'** to keep your current geometry.
\item User: confirm
\item Assistant: Your configuration is complete. The following items have been confirmed: shell clearance, nested core half length, nested core radius, and shell half length.
\end{itemize}
\textbf{Internal Chain (Natural Language):}
\begin{itemize}
\item Turn 1: action=summarize\_progress, phase=materials, complete=False.
\item Missing fields: materials.volume\_material\_map, source.type, source.particle, source.energy, source.position, source.direction, output.format, output.path, geometry.ask.shell.inner\_radius, geometry.ask.shell.thicknesses, geometry.ask.shell.half\_length, geometry.ask.shell.child\_radius, geometry.ask.shell.child\_half\_length, geometry.ask.shell.clearance.
\item LLM used=True; fallback=none; backend=llm\_slot\_frame+runtime\_semantic.
\item Turn 2: action=confirm\_overwrite, phase=source, complete=False.
\item Missing fields: none.
\item LLM used=True; fallback=none; backend=llm\_slot\_frame+runtime\_semantic.
\item Turn 3: action=finalize, phase=physics, complete=True.
\item Missing fields: none.
\item LLM used=True; fallback=none; backend=llm\_semantic\_frame+runtime\_semantic+confirmed\_pending\_overwrite.
\end{itemize}
\subsection{Case MTM18 (boolean\_cavity)}
\textbf{Dialogue:}
\begin{itemize}
\item User: 我想做一个挖空的盒体几何，但先只给你意图。
\item Assistant: 好的，我们已确认了您要创建“挖空盒体”的几何结构类型，并记录了盒体的尺寸。

接下来，为了完成这个挖空几何体的定义，还需要您提供构成布尔操作的两个实体（Solid A 和 Solid B）的具体尺寸信息。

**当前待补充信息：**
*   **Solid A** 的 X、Y、Z 尺寸。
*   **Solid B** 的 X、Y 尺寸。

请随时提供这些尺寸，我们继续配置。
\item User: 做减法布尔体：A 为 40 mm x 40 mm x 40 mm 盒体，B 为 20 mm x 20 mm x 20 mm 盒体，从 A 中减去 B；材料 ?；点源 gamma 0.3 MeV，位置 (0,0,-60) mm，方向 +z；物理列表 QBBC；输出 xml。
\item Assistant: 检测到您的新输入会覆盖已确认的几何参数。具体来说，布尔操作中的两个盒体尺寸将被更新：Solid A 将改为 40x40x40 mm，Solid B 将改为 20x20x20 mm。

如果您确认应用这些修改，请直接回复“确认”；如果您希望保留原来的尺寸值，请直接回复“保持原值”。
\item User: 确认
\item Assistant: 好的，配置已完成。当前已确认的配置项包括：布尔操作的两个盒体尺寸、输出路径和输出格式。
\end{itemize}
\textbf{Internal Chain (Natural Language):}
\begin{itemize}
\item Turn 1: action=summarize\_progress, phase=geometry, complete=False.
\item Missing fields: geometry.params.bool\_a\_x, geometry.params.bool\_a\_y, geometry.params.bool\_a\_z, geometry.params.bool\_b\_x, geometry.params.bool\_b\_y, geometry.params.bool\_b\_z, materials.volume\_material\_map, source.type, source.particle, source.energy, source.position, source.direction, physics.physics\_list, output.format, output.path, geometry.ask.boolean.solid\_a\_size, geometry.ask.boolean.solid\_b\_size.
\item LLM used=True; fallback=none; backend=llm\_slot\_frame+runtime\_semantic.
\item Turn 2: action=confirm\_overwrite, phase=materials, complete=False.
\item Missing fields: none.
\item LLM used=True; fallback=none; backend=llm\_slot\_frame+runtime\_semantic.
\item Turn 3: action=finalize, phase=source, complete=True.
\item Missing fields: none.
\item LLM used=True; fallback=none; backend=llm\_semantic\_frame+runtime\_semantic+confirmed\_pending\_overwrite.
\end{itemize}
\subsection{Case MTM19 (ring\_complete)}
\textbf{Dialogue:}
\begin{itemize}
\item User: Create a complete ring setup: 16 modules, each 6 mm x 6 mm x 2 mm, radius 32 mm, clearance 0.5 mm, material silicon; point gamma source 0.511 MeV at (0,0,0) mm; physics QBBC; output json.
\item Assistant: Your Geant4 configuration is complete. I have confirmed the volume-to-material binding, output path, clearance, and radius as specified.
\end{itemize}
\textbf{Internal Chain (Natural Language):}
\begin{itemize}
\item Turn 1: action=finalize, phase=materials, complete=True.
\item Missing fields: none.
\item LLM used=True; fallback=none; backend=llm\_slot\_frame+runtime\_semantic.
\end{itemize}
\subsection{Case MTM20 (grid\_complete)}
\textbf{Dialogue:}
\begin{itemize}
\item User: 请创建完整的阵列配置：3 x 3 盒体阵列，每个模块 12 mm x 12 mm x 3 mm，pitch\_x 15 mm，pitch\_y 15 mm，间隙 1 mm，材料 ???；点源 gamma 0.662 MeV，位置 (0,0,-90) mm，方向 +z；物理列表 QBBC；输出 root。
\item Assistant: 目前我已经同步了以下信息：
- **几何**：阵列结构、模块尺寸
- **源**：点源类型、伽马粒子、能量、位置和方向
- **物理**：物理列表
- **输出**：ROOT格式和路径

接下来，还需要您补充**每个几何体积对应的材料**，这样配置就完整了。
\end{itemize}
\textbf{Internal Chain (Natural Language):}
\begin{itemize}
\item Turn 1: action=summarize\_progress, phase=materials, complete=False.
\item Missing fields: materials.volume\_material\_map.
\item LLM used=True; fallback=none; backend=llm\_slot\_frame+runtime\_semantic.
\end{itemize}
\subsection{Case MTM21 (layered\_detector)}
\textbf{Dialogue:}
\begin{itemize}
\item User: I want a stacked detector slab, but the individual layer thicknesses are not fixed yet.
\item Assistant: Here's what we have so far:

✅ **Geometry:** You've chosen a stacked detector slab structure. We still need the footprint (size) of each layer and their individual thicknesses.

❌ **Materials:** We need to assign materials to your detector volumes.

❌ **Source:** We need the source type, particle type, and its energy.

❌ **Physics:** We need to select a physics list.

❌ **Output:** We need to choose an output format and a file path.

What would you like to set up next?
\item User: Use stack footprint 25 mm x 25 mm, thicknesses 1 mm, 2 mm, 1 mm, layer clearance 0.2 mm, outer box 30 mm x 30 mm x 12 mm, container clearance 0.5 mm, material silicon; beam electron 5 MeV from (0,0,-50) mm along +z; physics FTFP\_BERT; output json.
\item Assistant: Your new geometry settings would overwrite the previously confirmed stack and container dimensions.

**Pending Change:**
*   **Stack:** Footprint changes from \textasciitilde{}16.3 x 10.3 mm to 25 x 25 mm. Thicknesses change from \textasciitilde{}[6.7, 1.2, 5.6] mm to [1, 2, 1] mm. Layer clearance changes from \textasciitilde{}0.73 mm to 0.2 mm.
*   **Container Box:** Size changes from \textasciitilde{}18.1 x 45.5 x 16.7 mm to 30 x 30 x 12 mm. Container clearance changes from \textasciitilde{}0.87 mm to 0.5 mm.

Please reply **'confirm'** to apply these new values, or **'keep existing'** to retain the current geometry.
\item User: confirm
\item Assistant: Your configuration is complete. The following geometry items are confirmed: stack container clearance, stack container size, stack layer clearance, and stack layer thicknesses.
\end{itemize}
\textbf{Internal Chain (Natural Language):}
\begin{itemize}
\item Turn 1: action=summarize\_progress, phase=materials, complete=False.
\item Missing fields: materials.volume\_material\_map, source.type, source.particle, source.energy, source.position, source.direction, physics.physics\_list, output.format, output.path, geometry.ask.stack.footprint, geometry.ask.stack.thicknesses, geometry.ask.stack.layer\_clearance, geometry.ask.stack.parent\_size, geometry.ask.stack.nest\_clearance.
\item LLM used=True; fallback=none; backend=llm\_semantic\_frame+runtime\_semantic.
\item Turn 2: action=confirm\_overwrite, phase=source, complete=False.
\item Missing fields: none.
\item LLM used=True; fallback=none; backend=llm\_slot\_frame+runtime\_semantic.
\item Turn 3: action=finalize, phase=physics, complete=True.
\item Missing fields: none.
\item LLM used=True; fallback=none; backend=llm\_semantic\_frame+runtime\_semantic+confirmed\_pending\_overwrite.
\end{itemize}
\subsection{Case MTM22 (embedded\_sensor)}
\textbf{Dialogue:}
\begin{itemize}
\item User: 请创建完整的嵌套配置：外盒 80 mm x 80 mm x 80 mm，内嵌铅圆柱半径 15 mm、半长 25 mm，嵌套间隙 1 mm；材料 ?；点源 gamma 1 MeV，位置 (0,0,-100) mm，方向 +z；物理列表 Shielding；输出 csv。
\item Assistant: 配置已完成！当前已确认的设置包括：体积与材料绑定、输出路径、嵌套间隙（clearance）以及圆柱半径。
\end{itemize}
\textbf{Internal Chain (Natural Language):}
\begin{itemize}
\item Turn 1: action=finalize, phase=materials, complete=True.
\item Missing fields: none.
\item LLM used=True; fallback=none; backend=llm\_slot\_frame+runtime\_semantic.
\end{itemize}
\subsection{Case MTM23 (layered\_shield\_complete)}
\textbf{Dialogue:}
\begin{itemize}
\item User: Create a complete shell configuration: inner radius 12 mm, thicknesses 4 mm and 6 mm, half length 35 mm, nested core radius 8 mm, nested core half length 25 mm, clearance 1 mm, material iron; isotropic gamma source 0.8 MeV at (0,0,0) mm; physics QBBC; output hdf5.
\item Assistant: Your configuration is complete. The following items have been confirmed: box dimensions (x, y, z) and the assignment of materials to volumes.
\end{itemize}
\textbf{Internal Chain (Natural Language):}
\begin{itemize}
\item Turn 1: action=finalize, phase=materials, complete=True.
\item Missing fields: none.
\item LLM used=True; fallback=none; backend=llm\_slot\_frame+runtime\_semantic.
\end{itemize}
\subsection{Case MTM24 (boolean\_complete)}
\textbf{Dialogue:}
\begin{itemize}
\item User: 请创建完整布尔配置：做并集布尔体，A 为 30 mm x 30 mm x 20 mm 盒体，B 为 10 mm x 10 mm x 40 mm 盒体，与 A 做 union；材料 ?；束流 gamma 0.5 MeV，位置 (0,0,-70) mm，方向 +z；物理列表 QBBC；输出 json。
\item Assistant: 好的，目前配置进度如下：

**已完成设置：**
- **几何**：已确认布尔运算（并集）和两个盒体的尺寸。
- **源**：已设置伽马射线、能量 0.5 MeV、位置和方向。
- **物理**：已选定物理列表 QBBC。
- **输出**：已确定输出格式为 JSON 并指定了路径。

**仍需补充：**
- **材料**：需要您为几何体指定具体的材料。
\end{itemize}
\textbf{Internal Chain (Natural Language):}
\begin{itemize}
\item Turn 1: action=summarize\_progress, phase=materials, complete=False.
\item Missing fields: materials.volume\_material\_map.
\item LLM used=True; fallback=none; backend=llm\_slot\_frame+runtime\_semantic.
\end{itemize}
\end{document}